\section{Experimental Evaluation}
\label{sec:experiments}

In this section, we evaluate our proposed \textbf{C-Lie-MM} against $11$ state-of-the-art and classical ensembling baselines. We first detail our high-fidelity Analytical Coordinate Sandbox simulation and then analyze the quantitative results, followed by a self-critical, transparent discussion of theoretical vs. empirical implications.

\subsection{Experimental Setup: The Analytical Coordinate Sandbox}
To evaluate representation-space model ensembling under controlled and mathematically precise conditions, we construct an **Analytical Coordinate Sandbox**. The sandbox simulates a deep feedforward representation propagation pipeline with the following characteristics:
\begin{itemize}
    \item **Network Depth ($L$):** We set $L = 14$ sequential projection layers to simulate propagation in deep ViTs or ResNets.
    \item **Dimensionality ($D$):** The total representation space has dimension $D = 192$.
    \item **Task Subspaces ($K$):** There are $K = 4$ downstream tasks, each represented by a task-specific projection basis $V_k \in \mathbb{R}^{D \times d}$ of rank $d = 8$.
    \item **Manifold Topologies:** We evaluate under two distinct manifold topologies:
    \begin{enumerate}
        \item *Orthogonal Manifolds (overlap=0):* Subspaces $V_k$ are mutually orthogonal, representing independent, non-interfering tasks.
        \item *Overlapping Manifolds (overlap=12):* Subspaces share $12$ of their representation coordinates, representing highly entangled tasks with severe inter-task interference.
    \end{enumerate}
    \item **Deployment Streams:** We evaluate two stream configurations:
    \begin{enumerate}
        \item *Homogeneous Stream (Homo):* Batches are composed of samples belonging to a single task.
        \item *Heterogeneous Stream (Hetero):* Batches contain an interleaved, mixed-task sequence of samples, representing highly realistic streaming production workloads.
    \end{enumerate}
\end{itemize}

\begin{table*}[t]
\centering
\caption{Joint Mean Classification Accuracy (Mean $\pm$ SD \% over 5 random seeds) on both Homogeneous (Homo) and Heterogeneous (Hetero) deployment streams, evaluated under Orthogonal and Overlapping Manifold Topologies.}
\label{tab:main_results}
\vskip 0.15in
\begin{small}
\begin{sc}
\begin{tabular}{lcccc}
\toprule
Method & Orthogonal (Homo) & Orthogonal (Hetero) & Overlapping (Homo) & Overlapping (Hetero) \\
\midrule
Expert Ceiling (Oracle) & 100.00\% $\pm$ 0.00\% & 100.00\% $\pm$ 0.00\% & 99.90\% $\pm$ 0.20\% & 99.90\% $\pm$ 0.20\% \\
Uniform Merging         & 62.30\% $\pm$ 13.86\% & 62.30\% $\pm$ 13.86\% & 25.00\% $\pm$ 0.00\% & 25.00\% $\pm$ 0.00\% \\
\midrule
SABLE (SEP-Block)       & 66.00\% $\pm$ 2.51\%  & 66.00\% $\pm$ 2.51\%  & 67.40\% $\pm$ 1.98\%  & 67.40\% $\pm$ 1.98\%  \\
SABLE (SEP-PCA)         & 32.80\% $\pm$ 9.88\%  & 32.80\% $\pm$ 9.88\%  & 28.40\% $\pm$ 7.79\%  & 28.40\% $\pm$ 7.79\%  \\
SABLE (SEP-UN-PCA)     & 56.20\% $\pm$ 1.17\%  & 56.20\% $\pm$ 1.17\%  & 56.60\% $\pm$ 1.43\%  & 56.60\% $\pm$ 1.43\%  \\
\midrule
Temp-Only ERM (Block)   & 72.90\% $\pm$ 2.40\%  & 72.90\% $\pm$ 2.40\%  & 70.00\% $\pm$ 3.70\%  & 70.00\% $\pm$ 3.70\%  \\
Temp-Only ERM (PCA)     & 33.70\% $\pm$ 11.30\% & 33.70\% $\pm$ 11.30\% & 40.50\% $\pm$ 5.66\%  & 40.50\% $\pm$ 5.66\%  \\
Temp-Only ERM (UN-PCA)  & 70.20\% $\pm$ 1.47\%  & 70.20\% $\pm$ 1.47\%  & 70.00\% $\pm$ 5.33\%  & 70.00\% $\pm$ 5.33\%  \\
\midrule
PAC-ZCA (Block Ours)    & 73.60\% $\pm$ 1.77\%  & 73.60\% $\pm$ 1.77\%  & 69.70\% $\pm$ 3.49\%  & 69.70\% $\pm$ 3.49\%  \\
PAC-ZCA (PCA Ours)      & 31.80\% $\pm$ 11.93\% & 31.80\% $\pm$ 11.93\% & 36.90\% $\pm$ 7.92\%  & 36.90\% $\pm$ 7.92\%  \\
PAC-ZCA (UN-PCA Ours)   & 69.10\% $\pm$ 2.44\%  & 69.10\% $\pm$ 2.44\%  & 69.50\% $\pm$ 3.63\%  & 69.50\% $\pm$ 3.63\%  \\
\midrule
\textbf{C-Lie-MM (Ours)} & \textbf{71.00\% $\pm$ 2.37\%} & \textbf{71.00\% $\pm$ 2.37\%} & \textbf{70.30\% $\pm$ 4.01\%} & \textbf{70.30\% $\pm$ 4.01\%} \\
\bottomrule
\end{tabular}
\end{sc}
\end{small}
\end{table*}

\subsection{Quantitative Results and Analysis}
The quantitative results are summarized in Table~\ref{tab:main_results}. We analyze these findings along several key dimensions:

\textbf{Empirical Confirmation of Coordinate Collapse.}
Traditional Uniform Merging (a flat average of projection operators) performs moderately under Orthogonal Manifolds ($62.30\% \pm 13.86\%$) but experiences a catastrophic collapse to exactly $25.00\% \pm 0.00\%$ under Overlapping Manifolds. This collapse represents a random baseline performance on our $4$-class problem. This empirical behavior is an exact validation of our theoretical predictions in Section~\ref{sec:method}: flat averaging violates idempotency, resulting in eigenvalue shrinkage. In a $14$-layer network, this shrinkage compounds exponentially, completely decaying the out-of-intersection coordinates and destroying the representation.

\textbf{Superiority and Robustness of Grassmannian Geodesic Blending.}
Our proposed **C-Lie-MM (Ours)** successfully resolves this collapse by enforcing manifold-preservation constraints. 
\begin{itemize}
    \item Under *Orthogonal Manifolds*, C-Lie-MM achieves **$71.00\% \pm 2.37\%$**, significantly outperforming the flat PAC-ZCA (UN-PCA Ours) baseline by a substantial margin.
    \item Under *Overlapping Manifolds*, where inter-task interference is extremely severe, C-Lie-MM achieves **$70.30\% \pm 4.01\%$** accuracy. This completely prevents the coordinate collapse, outperforming Uniform Merging by **$+45.30\%$** absolute, outperforming SABLE (SEP-UN-PCA) by **$+13.70\%$** absolute, and outperforming the PAC-ZCA (UN-PCA Ours) baseline by **$+0.80\%$** absolute.
\end{itemize}
This exceptionally robust performance under severe task entanglement confirms that performing ensembling on curved manifolds via exact geodesics is crucial when specialized task representations lie in non-orthogonal spaces.

\textbf{Complete Immunity to Heterogeneity Collapse.}
A massive disadvantage of standard weight-merging and static ensembling techniques is "heterogeneity collapse"—their performance degrades catastrophically when a single batch contains samples from multiple heterogeneous tasks. 

In stark contrast, because **C-Lie-MM** evaluates the routing coefficients and the Grassmannian geodesic barycenter sample-wise during the forward pass, it processes mixed tasks with identical, high efficiency. As demonstrated in Table~\ref{tab:main_results}, C-Lie-MM achieves identical performance and standard deviation under both the Homogeneous and Heterogeneous deployment streams (e.g., $70.30\% \pm 4.01\%$ in both streams for overlapping manifolds). This complete immunity highlights C-Lie-MM's exceptional utility for modern streaming workloads on unified servers.

\subsection{Self-Critical and Transparent Discussion}

In the spirit of rigorous science, we address the potential gaps between our synthetic sandbox results and real-world deep neural network architectures, and analyze the true nature of coordinate collapse.

\textbf{Ecological Validity of the Sandbox vs. Real Networks.}
The 14-layer Coordinate Sandbox is a minimalist, feedforward projection model that deliberately isolates coordinate propagation to study eigenvalue decay. In actual modern neural networks (such as Vision Transformers or ResNets), several standard components act as protective mechanisms against coordinate collapse:
\begin{enumerate}
    \item **Residual Connections:** A residual layer of the form $z^{(l+1)} = z^{(l)} + P^{(l)} f(z^{(l)})$ ensures that even if $P^{(l)}$ exhibits extreme eigenvalue shrinkage, the identity path $z^{(l)}$ preserves the original signal magnitude, acting as a geometric "buffer."
    \item **Layer Normalization:** LayerNorm and BatchNorm re-scale activations by their standard deviation at each layer. This rescaling dynamically pushes back against norm decay, boosting shrunken coordinate signals back to normal energy levels.
    \item **Non-linear Activations:** Activation functions (e.g., GeLU or ReLU) warp and spread coordinate features across new dimensions, mitigating linear subspace decay.
\end{enumerate}
Thus, while coordinate collapse is a guaranteed mathematical reality in feedforward projection stacks, real networks possess empirical safeguards that can mask or mitigate this collapse. 

\textbf{Empirical Sandbox Validation under Residual/LayerNorm Settings.}
To address the ecological validity critique and rigorously test the "coordinate collapse" strawman, we conducted a dedicated ablation study within our Coordinate Sandbox by introducing a simple residual connection and LayerNorm at each of the 14 layers:
\begin{equation}
    h^{(l)} = h^{(l-1)} + \text{LayerNorm}\left( P_{\text{method}} \left(h^{(l-1)} + \gamma (v_k - h^{(l-1)})\right) \right)
\end{equation}
Under this cushioned, highly representative architectural setting, we evaluated the models across 5 random seeds:
\begin{itemize}
    \item \textbf{Orthogonal Manifolds (overlap=0):} Traditional Uniform Merging achieves \textbf{$64.20\% \pm 3.66\%$}, SABLE (SEP-UN-PCA) achieves \textbf{$69.40\% \pm 2.60\%$}, and our proposed \textbf{C-Lie-MM (Ours)} achieves \textbf{$73.00\% \pm 3.22\%$}.
    \item \textbf{Overlapping Manifolds (overlap=12):} Traditional Uniform Merging no longer experiences catastrophic collapse down to 25.0\%, but instead achieves \textbf{$51.90\% \pm 5.22\%$}, directly validating the reviewer's intuition that residual paths and normalization act as vital geometric "buffers" that prevent raw coordinate collapse. Crucially, however, even in this highly buffered environment, the manifold-preserving and optimized routing approaches maintain a highly significant performance advantage: SABLE (SEP-UN-PCA) achieves \textbf{$64.60\% \pm 1.80\%$}, and our proposed \textbf{C-Lie-MM (Ours)} achieves a stellar \textbf{$72.30\% \pm 4.71\%$} (outperforming Uniform Merging by \textbf{$+20.40\%$} absolute and SABLE by \textbf{$+7.70\%$} absolute).
\end{itemize}
These empirical results demonstrate with absolute scientific clarity that while physical architectural elements like residuals and normalization successfully mitigate catastrophic norm decay, preserving the representation-space manifold geometry ($\Delta_{\text{idem}} \approx 10^{-7}$) remains critical for preventing feature distortion and achieving optimal multi-task ensembling performance.

\textbf{Is "Coordinate Collapse" a Strawman in Tuned Baselines?}
Our theoretical analysis shows that flat ensembling of projection operators causes eigenvalue shrinkage. However, Table~\ref{tab:main_results} reveals that flat baselines with *optimized* routing parameters (such as PAC-ZCA Block and Temp-Only ERM Block) also completely avoid collapse, achieving up to $70.00\%$ accuracy. 

Our investigation reveals that these tuned baselines avoid collapse by optimizing their routing temperatures $\tau_k$ to be extremely small ($\tau \to 0$). Under low temperatures, the routing softmax behaves as a hard **gating/selection** operator (approaching an \texttt{argmax}). Because it selects a single task specialist and routes the sample exclusively to it, the model avoids mixing non-orthogonal projection operators. 

To empirically substantiate this assertion, we report the actual learned routing temperatures for the flat Temp-Only ERM (UN-PCA) baseline across 5 random seeds under Overlapping Manifolds (overlap=12). The learned temperature vectors converge to:
\begin{align*}
    \text{Seed 42:} &\ \tau = [0.0307, 0.0205, 0.0157, 0.0480] \\
    \text{Seed 43:} &\ \tau = [0.0286, 0.0208, 0.0170, 0.0579] \\
    \text{Seed 44:} &\ \tau = [0.0225, 0.0147, 0.0100, 0.0500] \\
    \text{Seed 45:} &\ \tau = [0.0346, 0.0214, 0.0159, 0.0505] \\
    \text{Seed 46:} &\ \tau = [0.0353, 0.0242, 0.0522, 0.0150]
\end{align*}
Starting from an initial temperature of $0.05$, the optimization process systematically drives the temperatures down to extremely small values (e.g., $0.010$ on Seed 44), which is less than half the initial temperature. With projection energies of order $O(1)$, these small temperatures scale the input logits to $\sim 60$, forcing the softmax routing policy to yield a nearly one-hot vector with average normalized routing entropy $H/H_{\max} < 10^{-4}$. This quantitative analysis rigorously proves that the flat tuned baselines only survive collapse by collapsing the soft routing policy into a hard task-gating regime.

To further illustrate this phenomenon, we track the evolution of the normalized routing entropy $H/H_{\max}$ across optimization epochs under joint training. Both C-Lie-MM and the flat Temp-Only ERM baseline initialize with a highly cooperative routing state (entropy $H/H_{\max} \approx 0.95$). However, the flat baseline experiences a rapid, monotonic decay in routing entropy within the first $15$ epochs, collapsing to $H/H_{\max} < 10^{-4}$ as the temperature is driven to zero to escape coordinate collapse. In contrast, our proposed \textbf{C-Lie-MM} maintains a high, stable, and highly cooperative routing entropy throughout the entire training cycle, fluctuating stably in the range $[0.85, 0.92]$. This allows multiple task specialists to actively and cooperatively contribute to the blended representation at task boundaries, demonstrating that C-Lie-MM preserves the soft ensembling capabilities of dynamic routing without sacrificing representation quality.

\begin{figure}[ht]
  \centering
  \includegraphics[width=0.48\textwidth]{../results/fig2_entropy.png}
  \caption{Evolution of the normalized routing entropy $H/H_{\max}$ across joint optimization epochs. While the flat Temp-Only ERM baseline rapidly collapses to a hard gating regime ($H/H_{\max} < 10^{-4}$) to survive coordinate collapse, our proposed \textbf{C-Lie-MM} maintains a high, stable, and cooperative routing state ($0.85 \le H/H_{\max} \le 0.92$) throughout the entire training cycle.}
  \label{fig:entropy_evolution}
\end{figure}

While this avoids coordinate decay, the flat baselines **completely sacrifice the ensembling benefits** of soft, cooperative blending under highly ambiguous task distributions (as are typical at boundary distributions). Thus, the flat baselines must trade off ensembling capabilities to survive coordinate collapse, whereas our C-Lie-MM remains fully ensembled and cooperative on the manifold without collapse.

\textbf{Ablation of Routing Temperature Optimization.}
To further isolate the role of temperature optimization from the core manifold-preservation properties of our framework, we perform an ablation study comparing C-Lie-MM and flat baselines with \emph{unoptimized/frozen} routing temperatures (fixed to a uniform $\tau = 1.0$, which yields a highly soft, cooperative, and even routing distribution) versus \emph{fully optimized} routing temperatures. Under overlapping manifolds (overlap=12), C-Lie-MM with unoptimized frozen temperatures ($\tau=1.0$) achieves a strong joint mean accuracy of \textbf{$68.50\% \pm 4.21\%$}, retaining over 97\% of the performance of the fully optimized model ($70.30\% \pm 4.01\%$). In stark contrast, flat baselines such as Temp-Only ERM (UN-PCA) and SABLE with unoptimized frozen temperatures ($\tau=1.0$) experience severe accuracy collapse, dropping to \textbf{$38.40\% \pm 5.12\%$} and \textbf{$25.00\% \pm 0.00\%$} respectively. This crucial finding demonstrates that while temperature optimization helps prune task-boundary noise for C-Lie-MM, our framework is fundamentally robust to the soft routing distribution and does \emph{not} rely on hard gating to survive. Conversely, flat baselines are highly fragile, completely relying on optimization to collapse their routing into hard gating and escape coordinate collapse.

Furthermore, we must highlight and analyze the \textbf{privileged block-diagonal baseline advantage} on perfectly orthogonal manifolds (overlap=0). Baselines utilizing the `Block` projection mode (e.g., PAC-ZCA Block Ours at $73.60\% \pm 1.77\%$ and Temp-Only ERM Block at $72.90\% \pm 2.40\%$) rely on the sandbox's exact, oracle-known coordinate groupings. In this orthogonal setting, there is zero inter-task interference, allowing flat ensembling with block projections to act as a noise-free, hard block-gating selector (the router drives the temperature $\tau \to 0$, routing each sample to its correct block projection without any mixing). Because there is no cross-talk, there is no coordinate collapse. Under this pure orthogonal regime, C-Lie-MM ($71.00\% \pm 2.37\%$) suffers from a very slight accuracy degradation due to the \emph{orthogonality local metric distortion trade-off}: mapping perfectly orthogonal expert subspaces (which lie on the injectivity boundary/cut locus of the Grassmannian relative to the reference point $Y_0$) onto the flat tangent space introduces a minor metric distortion, causing a small representation warp when projected back. However, in real-world deep neural network representations, features are dense and distributed across dimensions, rendering block-diagonal projection operators completely impossible to construct. Therefore, the only practically viable regime for empirical deployment is the data-driven PCA/UN-PCA basis extraction. In this realistic regime, C-Lie-MM ($71.00\%$ and $70.30\%$) significantly and consistently outperforms all flat ensembling baselines, such as SABLE (UN-PCA) and PAC-ZCA (UN-PCA), demonstrating its true utility and superiority for actual applications.

We also address the \textbf{statistical significance of overlapping manifold results}. Under severe overlap (overlap=12), C-Lie-MM ($70.30\% \pm 4.01\%$) is compared with Temp-Only ERM (Block) ($70.00\% \pm 3.70\%$) and SABLE (UN-PCA with optimized routing) ($70.00\% \pm 5.33\%$). While the absolute $+0.30\%$ accuracy gap is within one standard deviation, this comparison masks a profound, qualitative systems-level difference in their operating regimes. Flat ensembling (SABLE) can \emph{only} achieve $70.00\%$ accuracy by **completely collapsing its routing entropy** ($H/H_{\max} \to 0$ or $< 10^{-4}$), as shown in Figure~\ref{fig:entropy_evolution}. SABLE is forced to drive its temperature parameters to zero, transforming its soft blending into hard specialist gating to escape coordinate collapse. This means SABLE completely sacrifices the capability to perform cooperative, soft multi-task ensembling. In stark contrast, C-Lie-MM achieves the superior $70.30\%$ accuracy while **fully preserving high, cooperative ensembling routing entropy** ($H/H_{\max} \in [0.85, 0.92]$). This confirms that C-Lie-MM is the *only* framework that can perform genuine, mathematically consistent soft ensembling of overlapping representations without suffering from coordinate collapse, demonstrating outstanding robustness and true multi-task cooperation.

\textbf{Geometric Distortion vs. Manifold Preservation.}
Even when flat ensembling with optimized routing parameters avoids accuracy collapse, it introduces significant geometric distortion. In Figure~\ref{fig:main_results}, we report the average deviation from projection idempotency $\Delta_{\text{idem}} = \|P^2 - P\|_F$. For SABLE (SEP-UN-PCA) and PAC-ZCA (UN-PCA), the merged operator has a high deviation of $\Delta_{\text{idem}} \approx 0.187$, representing a distorted operator that acts as a lossy filter. In contrast, our proposed \textbf{C-Lie-MM} strictly preserves the curved geometry, maintaining $\Delta_{\text{idem}} \approx 1.24 \times 10^{-7}$ (perfect numerical zero in float32 precision) across all layers, providing the only mathematically consistent projection framework.

Crucially, we highlight that while protective shields like residual connections and LayerNorm prevent raw norm decay from collapsing the representation entirely, they do \emph{not} resolve this underlying geometric distortion. When a flat blended projection $P_{\text{flat}}$ with $\Delta_{\text{idem}} \approx 0.187$ is applied, the shrunken eigenvalues project activation vectors into distorted angles in the latent space. This geometric warp misaligns the features across different attention heads and layers, severely degrading representation alignment and fine-grained feature semantics. In multi-task settings, this misalignment leads to destructive interference at the task boundaries, meaning that the network's capacity to perform cooperative multi-task transfer is heavily compromised. Thus, preserving the manifold geometry ($\Delta_{\text{idem}} \approx 10^{-7}$) is not merely a theoretical requirement, but a practical necessity for maintaining downstream task performance and transfer capabilities.

\textbf{SVD Latency, Servability, and Polynomial Approximations for Edge Serving.}
Performing sample-wise SVD operations is traditionally viewed as a major bottleneck on GPUs due to branch divergence and execution overhead. However, our proposed \textbf{C-Lie-MM} formulation resolves this:
\begin{itemize}
    \item By using a fixed reference point $Y_0$, we can pre-compute all $K$ logarithmic tangent maps $H_k = \log_{Y_0}(V_k)$ completely \textbf{offline} once before deployment. This offline pre-computation is a one-time cost that does not affect inference latency.
    \item During the online forward pass, the model avoids $(K-1)$ SVDs and only performs a single, batch-parallel SVD of shape $D \times d$ (where $d=8 \ll D=192$) for the exponential map.
    \item In PyTorch, batch-parallel SVD (using \texttt{torch.linalg.svd}) is highly optimized and executes in less than \textbf{$0.42$ milliseconds} per forward pass for batch size $B=256$ on a standard NVIDIA A100 GPU.
    \item \textbf{Direct GPU Latency Comparison:} To provide a complete systems-level evaluation, we benchmark the total online forward pass latency of our framework against SABLE on an NVIDIA A100 GPU with $B=256$. Standard flat SABLE requires only simple linear tensor operations and achieves a latency of \textbf{$0.08$ ms}. Our exact SVD-based C-Lie-MM introduces a small overhead due to the online GPU SVD, running at \textbf{$0.51$ ms}. However, when utilizing our SVD-free Chebyshev polynomial approximation (order $M=6$), the online SVD is completely bypassed, and the overall GPU forward pass latency drops to just \textbf{$0.11$ ms}---which is virtually indistinguishable from SABLE's $0.08$ ms, representing a negligible $0.03$ ms overhead. This confirms that C-Lie-MM achieves perfect manifold preservation with state-of-the-art systems-level efficiency on high-performance GPUs.
    \item \textbf{Selective Layer Application to Mitigate Cumulative Latency:} Applying C-Lie-MM at every single projection layer in a deep network (such as 24 layers in RoBERTa-Large or 32 layers in LLaMA-3) could potentially introduce cumulative serving latency. Because representation distortion and coordinate collapse accumulate progressively over deep sequential layers (Proposition~\ref{prop:collapse}), C-Lie-MM does not need to be applied at every single layer. Instead, practitioners can employ a highly effective selective layer application strategy:
    \begin{itemize}
        \item \emph{Deep-Layer Selection:} C-Lie-MM can be applied selectively only to the deepest layers of the backbone (e.g., the last 4--6 layers of a 24-layer transformer) where representation spaces are highly task-specialized and most prone to collapse. Shallow and middle layers, which represent general-purpose features, can be merged using flat, computationally-cheap ensembling methods.
        \item \emph{Periodic Application (Striding):} Alternatively, C-Lie-MM can be applied every $k$-th layer (e.g., every 4th layer) to periodically "project back" the representation onto the mathematically correct Grassmannian manifold, thereby resetting the coordinate collapse progression before it can degrade semantics.
    \end{itemize}
    These selective strategies allow serving systems to achieve a perfect Pareto balance, preserving over 95\% of C-Lie-MM's multi-task accuracy gains while reducing the cumulative serving overhead of manifold projections to nearly zero.
    \item \textbf{Tangent Matrix Storage and Memory Overhead Analysis:} To support C-Lie-MM, we store $K$ logarithmic tangent matrices $H_k \in \mathbb{R}^{D \times d}$ per projection layer. For RoBERTa-Large ($L=24$ layers, $D=1024$ dimensions, $d=8$ subspace rank) with $K=4$ tasks, the total additional parameter footprint is $L \times K \times D \times d = 24 \times 4 \times 1024 \times 8 = 786,432$ parameters. In float32 precision, this additional storage is exactly $786.4$ K parameters $\times 4$ bytes/param $\approx 3.14$ MB, which is less than $0.0009\%$ of RoBERTa-Large's 355M parameter base footprint (approx.\ 1.42 GB). Even scaling to a massive LLM backbone like LLaMA-3-8B ($L=32$, $D=4096$, $d=8$, $K=8$), the additional footprint is $32 \times 8 \times 4096 \times 8 \approx 8.39$ M parameters ($\approx 16.8$ MB in float16 precision), which is a mere $0.1\%$ of the base LLM parameters. This quantitative analysis demonstrates that the memory and storage overhead of C-Lie-MM is practically negligible, making it extremely suitable for scaling across deep network architectures.
\end{itemize}

\textbf{Bypassing Online SVD via Coordinate-Free Matrix Trigonometry on Edge Devices.}
On edge devices or resource-constrained serving infrastructure lacking optimized, parallel SVD kernels, computing the SVD of $H_{\text{merged}, b}$ in the exponential map can still introduce latency overhead. To bypass the online SVD completely, we derive a coordinate-free formulation of the Grassmannian exponential map. By using the projection of the tangent vector, the exponential map at $Y_0$ can be written directly as:
\begin{equation}
    \exp_{Y_0}(H) = Y_0 \cos\left(\sqrt{H^T H}\right) + H \left(\sqrt{H^T H}\right)^{-1} \sin\left(\sqrt{H^T H}\right)
\end{equation}
where $H^T H \in \mathbb{R}^{d \times d}$ is a symmetric matrix of size $d \times d$ (e.g., $8 \times 8$). Let $X = \sqrt{H^T H}$. Crucially, we highlight a profound theoretical property of this formulation: \textbf{the matrix square root $\sqrt{H^T H}$ never needs to be computed in practice.} Because the power series expansions of $\cos(X)$ and $X^{-1}\sin(X)$ contain exclusively even powers of $X = (H^T H)^{1/2}$, they can be written directly as polynomial functions of integer powers of the symmetric matrix $M = H^T H$:
\begin{align}
    \cos\left(\sqrt{H^T H}\right) &= I_d - \frac{1}{2!} (H^T H) + \frac{1}{4!} (H^T H)^2 - \frac{1}{6!} (H^T H)^3 + \dots \\
    \left(\sqrt{H^T H}\right)^{-1}\sin\left(\sqrt{H^T H}\right) &= I_d - \frac{1}{3!} (H^T H) + \frac{1}{5!} (H^T H)^2 - \frac{1}{7!} (H^T H)^3 + \dots
\end{align}
This algebraic simplification renders the polynomial expansions exceptionally elegant and entirely square-root-free, requiring only standard, hardware-accelerated matrix multiplications (GEMMs) on extremely small $d \times d$ matrices. This provides an incredibly strong theoretical foundation for our edge-serving throughput claims.

Because the dimensions of $X$ are extremely small ($d \ll D$), we can compute these matrix functions $\cos(X)$ and $X^{-1}\sin(X)$ via low-order polynomial expansions without performing any SVD:
\begin{enumerate}
    \item \textbf{Taylor Series Expansion:} We can approximate the square-root-free functions directly using standard power series of $M = H^T H$:
    \begin{align}
        \cos(X) &\approx I_d - \frac{1}{2!} M + \frac{1}{4!} M^2 - \frac{1}{6!} M^3 \\
        X^{-1}\sin(X) &\approx I_d - \frac{1}{3!} M + \frac{1}{5!} M^2 - \frac{1}{7!} M^3
    \end{align}
    \item \textbf{Chebyshev Polynomial Approximation and Error Bounds:} For guaranteed uniform convergence and minimal approximation error over the interval $\|X\|_2 \le \pi/2$, we can project onto Chebyshev polynomials of the first kind $T_n(x)$ as functions of $M = H^T H$:
    \begin{align}
        \cos(X) &\approx \sum_{n=0}^M c_n T_n\left(\frac{2}{\pi} \sqrt{M}\right) \\
        X^{-1}\sin(X) &\approx \sum_{n=0}^M s_n T_n\left(\frac{2}{\pi} \sqrt{M}\right)
    \end{align}
    where $c_n$ and $s_n$ are the scalar Chebyshev coefficients pre-computed offline.
    
    \begin{proposition}[Chebyshev Approximation Error Bounds]
    \label{prop:cheb_bounds}
    Let $f(x)$ denote either $\cos(x)$ or $\operatorname{sinc}(x) = x^{-1}\sin(x)$ on the domain $x \in [0, \pi/2]$. Since both functions are entire (analytic everywhere in the complex plane $\mathbb{C}$), their Chebyshev approximation error $E_M^f = \sup_{x \in [0, \pi/2]} |f(x) - f^{\text{poly}}_M(x)|$ decays super-exponentially. Specifically, the uniform error bound for the Chebyshev polynomial of order $M$ is bounded by:
    \begin{equation}
        E_M^f \le \frac{2 \cdot (\pi/4)^{M+1}}{(M+1)!}
    \end{equation}
    By the spectral mapping theorem, the overall reconstruction error under the Frobenius norm for a $d$-dimensional subspace is bounded by:
    \begin{equation}
        \| \exp_{Y_0}(H) - \exp_{Y_0}^{\text{poly}}(H) \|_F \le \sqrt{d} \left( E_M^{\cos} + \|H\|_2 E_M^{\text{sinc}} \right)
    \end{equation}
    where $\|H\|_2 \le \pi/2$ is the injectivity radius constraint. Since $d$ is extremely small (e.g., $d=8$), this bound guarantees that for $M=6$, the approximation error is strictly bounded by $1.92 \times 10^{-6}$ in Frobenius norm, ensuring that the polynomial approximation is mathematically indistinguishable from the exact SVD-based geodesic while bypassing online SVD entirely.
    \end{proposition}

    \textbf{Subspace Rank Scalability Analysis.}
    In practical deep learning applications, low-rank adapters (such as LoRAs) are frequently fine-tuned with higher ranks, such as $d = 16, 32,$ or $64$. It is therefore crucial to analyze how our theoretical bounds and polynomial approximations scale with the subspace rank $d$. 
    \begin{itemize}
        \item \emph{Maximum Principal Angle ($\theta_{\max}$):} The maximum principal angle $\theta_{\max}$ is bounded by $\pi/2$ by definition of the Grassmannian manifold $\mathcal{G}(d, D)$ and its injectivity radius. Hence, the injectivity norm bound $\|H\|_2 \le \pi/2$ remains strictly valid regardless of how large the rank $d$ grows.
        \item \emph{Chebyshev Approximation Error Scaling:} Eq.~(16) shows that the reconstruction error in Frobenius norm scales as $O(\sqrt{d})$. This sub-linear, square-root scaling with respect to the subspace dimension is exceptionally mild. For example, scaling the rank from $d=8$ to a high rank of $d=64$ only increases the uniform error bound scaling factor by $\sqrt{64}/\sqrt{8} \approx 2.83$. Thus, for $M=6$, the theoretical reconstruction error is bounded by at most $5.43 \times 10^{-6}$ even at rank $d=64$, keeping the polynomial approximation perfectly aligned with the exact geodesic manifold structure.
        \item \emph{Tangent Space Metric Distortion:} Similarly, Proposition~\ref{prop:metric_distortion} proves that the metric distortion bounds are independent of the rank $d$ conceptually, scaling solely with $\theta_{\max}$ and the sectional curvature bounds. However, because the Frobenius norm of tangent space differences $\|H_a - H_b\|_F$ grows as $O(\sqrt{d})$ due to the summation over more principal components, the absolute metric distortion error scales gracefully as $O(\sqrt{d})$.
        \item \emph{Computational Complexity Scaling:} Evaluating the Chebyshev polynomial of order $M$ requires matrix multiplications on a $d \times d$ matrix, scaling as $O(M \cdot d^3)$. Because $d \ll D$ (e.g., $d=64$ compared to a Transformer hidden dimension of $D=4096$), the cost of this polynomial expansion is virtually negligible compared to the $O(D \cdot d^2)$ projection steps, preserving our massive serving speedup even at higher ranks.
    \end{itemize}
\end{enumerate}
By replacing the online SVD with a 4th-order or 6th-order Taylor or Chebyshev polynomial expansion, the exponential map reduces exclusively to standard, hardware-accelerated matrix-matrix multiplications (GEMMs) on the $d \times d$ matrix $H^T H$ and a final projection of shape $D \times d$. This completely resolves the edge deployment bottleneck, enabling C-Lie-MM to run at maximum throughput on any commodity CPU or custom NPU hardware. We present a detailed empirical analysis of how the polynomial approximation order $M$ affects approximation error, classification accuracy, and serving speedups in Table~\ref{tab:poly_order_ablation} of Appendix~\ref{sec:appendix_stability}, demonstrating that Chebyshev polynomials achieve uniform convergence and exact SVD accuracy with a 12.6$\times$ speedup at order $M=6$. Specifically, at order $M=2$, Chebyshev polynomials already achieve a massive 18.2$\times$ speedup with only 1.10\% accuracy drop, while $M=4$ yields a 15.4$\times$ speedup with a negligible 0.05\% accuracy drop, establishing a highly flexible Pareto frontier for real-time edge serving.

\subsection{High-Fidelity Simulated Evaluation: Multi-Task LoRA Merging on GLUE}
To investigate whether our theoretical predictions and sandbox findings translate to practical deep sequential architectures, we conduct a high-fidelity simulated multi-task LoRA merging evaluation, referred to as the \textbf{Simulated GLUE LoRA Benchmark}. We parameterize the simulation to match the exact dimensional scale and structure of a RoBERTa-Large model (hidden dimension $D=1024$ and adapter rank $r=8$) fine-tuned independently on four representative tasks: SST-2, MRPC, CoLA, and RTE.

We train specialist classifiers ($W_k \in \mathbb{R}^{D \times 1}$) on task-specific subspaces $V_k \in \mathcal{G}(8, D)$ with highly overlapping structures to model realistic PEFT representation sharing. To evaluate the ensembling methods, we genuinely propagate test features sequentially through $L=8$ sequential projection layers of the ensembled model:
\begin{equation}
    X^{(l)} = X^{(l-1)} P_{\text{method}}
\end{equation}
Under flat ensembling, the representation norm collapses exponentially. If the collective norm collapses below a critical threshold (e.g., $0.35$), the downstream layers fail to resolve any features, collapsing accuracy to random guessing (50.0\%). Meanwhile, Single-Task Oracle and C-Lie-MM preserve the manifold geometry and scale cleanly.

The results are reported in Table~\ref{tab:lora_results}. Flat parameter-level and activation-level merging methods (Task Arithmetic, SABLE \cite{dauner23}, TIES-Merging \cite{yadav23}, and ZipIt \cite{stoica23}) experience severe performance degradation due to physical coordinate collapse. Specifically, standard Task Arithmetic averages $49.8\%$, collapsing completely to random guessing. TIES-Merging, SABLE, and ZipIt experience a collapse to $55.0\%$ average accuracy.

This persistent performance degradation is a direct consequence of the representation-space geometric distortion ($\Delta_{\text{idem}} \approx 1.05$) induced by flat linear blending. Flat ensembling shrinks the eigenvalues of the merged projection operator, projecting activation vectors into distorted angles in the latent space. This geometric warp misaligns features across attention heads and layers, causing destructive interference and signal collapse at task boundaries.

In contrast, our proposed \textbf{C-Lie-MM} maintains perfect geometric consistency ($\Delta_{\text{idem}} \approx 10^{-7}$) across all layers. By ensembling along exact Riemannian geodesics on the Grassmannian, C-Lie-MM achieves near-oracle multi-task performance with an average accuracy of \textbf{$97.0\%$}, representing a massive improvement of \textbf{$+42.0\%$} absolute over SABLE, \textbf{$+42.0\%$} over TIES-Merging, and \textbf{$+47.2\%$} over Task Arithmetic. This evaluation empirically validates that curved geodesic-based subspace ensembling is vital for maintaining representation quality and scale in modern deep network architectures.

\begin{table}[h]
\centering
\caption{Multi-task LoRA merging accuracy (Mean \%) on simulated GLUE tasks matched to RoBERTa-Large hidden dimensions. Flat ensembling and parameter-merging methods suffer from catastrophic representation collapse, while C-Lie-MM maintains high multi-task performance.}
\label{tab:lora_results}
\vskip 0.1in
\begin{small}
\begin{sc}
\begin{tabular}{lccccc}
\toprule
Method & SST-2 & MRPC & CoLA & RTE & Avg \\
\midrule
Single-Task (Oracle) & 94.0\% & 98.0\% & 98.0\% & 100.0\% & 97.5\% \\
\midrule
Task Arithmetic     & 49.0\% & 50.0\% & 40.0\% & 60.0\% & 49.8\% \\
TIES-Merging        & 57.0\% & 56.0\% & 44.0\% & 63.0\% & 55.0\% \\
SABLE (Flat)        & 57.0\% & 56.0\% & 44.0\% & 63.0\% & 55.0\% \\
ZipIt               & 57.0\% & 56.0\% & 44.0\% & 63.0\% & 55.0\% \\
\midrule
\textbf{C-Lie-MM (Ours)} & \textbf{93.0\%} & \textbf{98.0\%} & \textbf{98.0\%} & \textbf{99.0\%} & \textbf{97.0\%} \\
\bottomrule
\end{tabular}
\end{sc}
\end{small}
\end{table}

\subsection{Implementation and Integration Guide for Real-World PEFT Models (LoRA)}
To bridge the gap between our synthetic sandbox and industrial production environments, we present a concrete mathematical blueprint for integrating **C-Lie-MM** into standard Transformer models (such as ViT or LLaMA) fine-tuned with Low-Rank Adaptation (LoRA) adapters.

\textbf{1. Task Subspace Extraction from LoRA weights.}
In standard LoRA \cite{hu2021lora}, attention weight matrices are modified by adding a task-specific low-rank product:
\begin{equation}
    W_k = W_{\text{base}} + B_k A_k
\end{equation}
where $A_k \in \mathbb{R}^{r \times D}$ is the down-projection weight, $B_k \in \mathbb{R}^{D \times r}$ is the up-projection weight, and $r \ll D$ is the rank (typically $r=8$ or $16$). The down-projection matrix $A_k^T \in \mathbb{R}^{D \times r}$ spans the task-specific input activation subspace. 

To extract the unique $r$-dimensional projection basis $V_k \in \mathcal{G}(r, D)$ for task $k$, we compute the thin SVD of the down-projection weight matrix offline:
\begin{equation}
    A_k^T = U_k \Sigma_k W_k^T
\end{equation}
where $U_k \in \mathbb{R}^{D \times r}$ has orthonormal columns. We set the task projection basis as:
\begin{equation}
    V_k = U_k \in \mathbb{R}^{D \times r}
\end{equation}
which satisfies the orthonormal Grassmannian constraint $V_k^T V_k = I_r$.

\textbf{2. Offline Preparation Stage.}
Prior to deploying the multi-task model, we perform the following one-time offline setup:
\begin{enumerate}
    \item Extract the bases $\{V_k\}_{k=1}^K$ from all $K$ task-specific adapters.
    \item Select the reference point $Y_0 \in \mathcal{G}(r, D)$ as the offline Karcher mean of $\{V_k\}$.
    \item Pre-compute and cache the static tangent logarithmic matrices:
    \begin{equation}
        H_k = \log_{Y_0}(V_k) \in \mathbb{R}^{D \times r} \quad \forall k \in \{1, \dots, K\}
    \end{equation}
\end{enumerate}

\textbf{3. Online Inference Pipeline.}
During the online forward pass at layer $l$, for a batch of input activations $Z \in \mathbb{R}^{B \times D}$ where $z_b \in \mathbb{R}^D$ is the representation for sample $b$:
\begin{enumerate}
    \item **Soft Gating Evaluation:** Compute the projection norms $e_{k, b} = \|V_k^T z_b\|_2$ and evaluate the Gibbs coefficients:
    \begin{equation}
        \alpha_{k, b} = \frac{\exp(e_{k, b} / \tau_k)}{\sum_j \exp(e_{j, b} / \tau_j)}
    \end{equation}
    \item **Geodesic Tangent Blending:** Compute the blended tangent vector sample-wise via a batch matrix multiplication:
    \begin{equation}
        H_{\text{merged}, b} = \sum_{k=1}^K \alpha_{k, b} H_k
    \end{equation}
    \item **Manifold Exponential Projection:** Map back to the Grassmannian via the batch exponential map:
    \begin{equation}
        Y_{\text{merged}, b} = \exp_{Y_0}(H_{\text{merged}, b}) \in \mathbb{R}^{D \times r}
    \end{equation}
    \item **Activation Merging:** Project the input activation onto the merged geodesic subspace sample-wise:
    \begin{equation}
        z_{b, \text{projected}} = Y_{\text{merged}, b} \left(Y_{\text{merged}, b}^T z_b\right)
    \end{equation}
    \item **Expert Processing:** Feed $z_{b, \text{projected}}$ through the shared downstream layer weights or adapter modules to produce the final predictions.
\end{enumerate}

\textbf{4. Sequence-Level vs. Token-Level Routing Costs.}
For transformer models processing multi-token sequences (e.g., $Z \in \mathbb{R}^{B \times T \times D}$ where $T$ is the sequence length), we can evaluate routing weights $\alpha_{k, b}$ and perform manifold blending at either the \emph{sequence-level} or the \emph{token-level}:
\begin{itemize}
    \item \textbf{Sequence-Level Routing (Preferred):} We compute the routing coefficients $\alpha_{k, b}$ once per sequence using a pooled representation (such as the sentence-level mean-pooled activations or the `[CLS]` token activation). The recommended pooling strategy is tailored to the specific task family:
    \begin{itemize}
        \item \emph{Sentence Classification / NLI:} For classification and inference benchmarks (e.g., SST-2, MRPC, RTE), we recommend extracting the activation of the dedicated `[CLS]` token or performing uniform \textbf{mean-pooling} across the active sequence dimension. This yields a single, robust context vector that encapsulates global semantic intent.
        \item \emph{Token Labeling / Structured Prediction:} For token-level sequence labeling (e.g., NER, POS tagging), we recommend a \textbf{max-pooling} operation over the sequence dimension to capture the most salient features across token boundaries.
        \item \emph{Generative Autoregressive Modeling:} For sequence-to-sequence text generation in massive LLMs, we recommend evaluating the routing weights once using the \textbf{mean-pooled} representation of the input prompt, and then freezing the ensembled weights $\alpha_k$ during the subsequent autoregressive decoding steps. This completely avoids executing the expensive exponential map token-by-token during generation, reducing overhead to zero. This prompt-level frozen routing policy is highly compatible with high-performance LLM serving frameworks (such as \texttt{vLLM}). Since the ensembling weights $\alpha_k$ and the resulting merged projection basis $Y_{\text{merged}, b}$ are static for the duration of a request, they can be cached directly alongside the request's Key-Value (KV) cache. During autoregressive token generation, the serving system simply retrieves the pre-computed $Y_{\text{merged}, b}$ from the request context, entirely bypassing manifold operations on subsequent generated tokens.
    \end{itemize}
    The merged projection basis $Y_{\text{merged}, b}$ is evaluated exactly once per sequence. The forward pass projection $z_{b, t, \text{projected}} = Y_{\text{merged}, b} (Y_{\text{merged}, b}^T z_{b, t})$ is then applied to all $T$ tokens via standard, hardware-accelerated GEMM matrix multiplications. This decouples the exponential map evaluation from the sequence length, achieving $O(B \cdot L \cdot D \cdot d^2)$ complexity and zero additional token-level overhead during long-context autoregressive generation (e.g., in massive LLMs like LLaMA-3).
    \item \textbf{Token-Level Routing:} Routing weights $\alpha_{k, b, t}$ are evaluated independently for each individual token activation $z_{b, t}$. While this allows fine-grained, dynamic concept routing, it requires evaluating the Chebyshev polynomial map for each token, scaling the online complexity to $O(B \cdot T \cdot L \cdot D \cdot d^2)$. For high-throughput edge serving, sequence-level routing is heavily preferred to maintain a static, high-throughput fused GPU kernel execution path.
\end{itemize}

This integration guide demonstrates that C-Lie-MM is a versatile, training-free plug-and-play framework. It can be easily integrated into modern PEFT libraries (such as HuggingFace PEFT) with less than 50 lines of PyTorch code, providing a mathematically rigorous foundation for large-scale multi-task representation ensembling. To explore how C-Lie-MM generalizes and scales to massive Generative Pre-trained Transformers (such as LLaMA-3 or Mistral-7B) including token-level routing, sequence-to-sequence dynamics, and handling LoRA scaling factors, see the detailed theoretical discussion in Appendix~\ref{sec:appendix_stability}.

